\section{Threats to Validity and Scope of the Claim}
\label{sec:threats}

We state the boundaries of the evidence above explicitly.

\paragraph{We report \(\Sigma\)TTFT, not end-to-end latency.} Summed TTFT is the
quantity the mechanism acts on, and it is what we measure. We do not claim a
proportional end-to-end win: the fraction of a request's wall-clock time that
\(\Sigma\)TTFT represents depends on decode lengths, tool latency, and retrieval
cost, none of which \sysname{} changes.

\paragraph{The Oracle column is an upper bound, and should be read as one.} Oracle
is measured by replaying an identically-sized prompt and reading the second serve,
so it reflects what the hardware does when every token is already cached. It
upper-bounds what \emph{any} KV-reuse system could achieve in this setting. The
72\% figure in Table~\ref{tab:percomponent} is a fraction of that bound, not a
fraction of total request latency.

\paragraph{The hardware constrains the mechanism.} On an RTX 3090, a 4B model is
large relative to the device. Prefill cannot be overlapped with decode without
harming decode latency, so prefill--decode aggregation is disabled in these runs.
This cuts both ways: it removes a source of interference that would complicate the
measurement, but it also means the spatial-idle path of Section~\ref{sec:design}
is not exercised here, and results on hardware where aggregation is viable may
differ in both directions.

\paragraph{The baseline is vLLM's APC.} We build on vLLM, so the baseline adopts
vLLM's automatic prefix caching. We do not compare against the workflow-aware KV
management systems discussed in Section~\ref{sec:related}; those systems govern a
different stage of the KV lifecycle, and a fair comparison would require porting
them onto the same runtime and hardware.

\paragraph{Two topologies, one model, one GPU.} The evidence is two case studies
at a single model size on a single device, with 20 requests each. They were chosen
to be structurally different --- a static fan-out with a supervisor branch, and a
data-dependent loop with a self-referential carry --- but they do not establish
how the approach behaves under concurrency, at other model sizes, or on workloads
whose idle windows are shorter than the prefill work available to fill them. The
broader evaluation is ongoing.

\paragraph{A configuration caveat, not a result caveat.} The measurements use
compiler-derived speculative prefill with no probe decode and no deploy-time
prewarm. The zero-gain rows for \texttt{route} and the four \texttt{analyze}
stages are therefore properties of this configuration. We report them as measured
and do not extrapolate what the probe mechanism would recover.
